\section{Streaming}
\subsection{Protocolli real-time (push) e streaming (pull)}
Esistono diverse classi di applicazioni con requisiti differenti che devono operare tutti sulla stessa infrastruttura di
rete IP già esistente. È necessario sviluppare protocolli di trasmissione appositi per soddisfare i requisiti di ciascuna
classe sacrificando alcune funzionalità non necessarie.

\subsubsection*{Protocollo Real Time - push}
Il protocollo real time o push è un protocollo leggero ed essenziale pensato per la trasmissione dati di comunicazioni
real-time con requisiti di latenza molto stringenti. È usato per servizi di videoconferenza, cloud gaming e chirurgia
remota. È composto da:
\begin{itemize}
	\item \textbf{UDP} o User Datagram Protocol: protocollo di trasporto che offre solo un servizio di consegna dei pacchetti
	senza garanzia di consegna, è estremamente leggero e veloce, ma molto essenziale
	\item \textbf{RTP} o Real-Time Transport Protocol: aggiunge ai pacchetti UDP un header con informazioni di sequenza e
	timestamp per sincronizzazione, ordinamento e gestione del jitter
	\item \textbf{WebRTC} o Web Real-Time Communication: framework di sistema, aggiunge cifratura dei dati, meccanismi per
	attraversare NAT e firewall, accesso hardware per gestire microfono e webcam direttamente dal browser
\end{itemize}

\subsubsection*{Protocollo Streaming - pull}
Il protocollo streaming o pull è un protocollo più complesso e pesante pensato per la trasmissione dati di contenuti
multimediali che richiedono maggiore affidabilità della trasmissione, ma tollerano una latenza maggiore. È composto da:
\begin{itemize}
	\item \textbf{CDN} o Content Delivery Network: rete di server distribuiti geograficamente che memorizzano copie dei
	media da trasmettere, permettono di ridurre la latenza siccome avvicinano i contenuti agli utenti finali, riducono
	il traffico sulla rete principale siccome la distribuzione avviene tramite reti locali e aumentano l'affidabilità
	grazie alla ridondanza dei contenuti
	\item \textbf{QUIC/TCP} o Quick UDP Internet Connections/Transmission Control Protocol: protocolli di trasporto che
	offrono un servizio di trasmissione a pacchetti affidabile, garantendone la consegna
	\item \textbf{HTTP/3} o Hypertext Transfer Protocol: protocollo di livello applicativo che permette la trasmissione
	di contenuti multimediali tramite il protocollo di trasporto QUIC/TCP
	\item \textbf{MPEG-DASH/HLS} o Dynamic Adaptive Streaming over HTTP/HTTP Live Streaming: protocolli di streaming che
	gestiscono la trasmissione dei contenuti multimediali tramite il protocollo HTTP/3, adattando la qualità in base alle
	prestazioni della connessione dell'utente finale
	\item \textbf{codec video}: algoritmi di compressione e decompressione che avvengono all'interno dell'applicativo nella
	macchina client
\end{itemize}

\subsection{Adattamento del tasso di codifica video}
\subsubsection*{Scelta del tasso di codifica video}
In generale il tasso di codifica video deve essere inferiore al throughput della rete \(R_C \leq S\). Ad oggi si
conoscono strategie avanzate per adattare il tasso di codifica video in tempo reale alle prestazioni della rete,
ma spesso le prestazioni della rete non si conoscono a priori, specialmente se le fasi di codifica e di trasmissione
avvengono in momenti diversi o in contesti di distribuzione in cui un unico server offre servizi a molti client.

Se il tasso di codifica è molto più basso rispetto al throughput disponibile, la qualità video sarà molto bassa e
l'esperienza utente sarà scadente. Se il tasso di codifica è più alto rispetto al throughput disponibile, si verificheranno
fenomeni di buffer overflow e di perdita dei pacchetti, causando interruzioni nella riproduzione del video e un'esperienza
utente scadente.

Di seguito si analizzano alcune strategie per garantire la massima qualità video possibile limitatamente alle prestazioni
della rete, gestendo adeguatamente il tasso di codifica video.

\subsubsection*{Compressione estrema}
L'approccio più semplice è quello di comprimere il video al massimo, riducendo il bitrate al minimo possibile. In questo
modo si garantisce un bitrate inferiore al throughput della rete, ma la qualità video sarà molto bassa anche quando la rete
permette la trasmissione di video con qualità maggiore a bitrate più elevati. L'esperienza utente sarà quindi scadente.

\subsubsection*{Transcodifica}
L'approccio prevede che le macchine router o server effettuino una transcodifica del video in tempo reale per abbassare
la qualità in caso di congestione della rete. Questo approccio ha alcune criticità fondamentali:
\begin{itemize}
	\item il costo computazionale è insostenibile per le macchine router o server
	\item non è possibile accedere ai dati trasmessi in quanto cifrati end-to-end
	\item un banale packet drop sistematico non può funzionare in quanto gli algoritmi sono altamente predittivi e poco
	tolleranti ad errori sui frame
\end{itemize}

\subsubsection*{Scalable video coding - SVC}
L'approccio consiste nell'utilizzare dal lato server una codifica video scalabile, ovvero composta da più flussi video
di qualità sempre crescente. Si codifica inizialmente un video in un flusso base con alti tassi di compressione e quindi
bassa qualità video e basso bitrate, in modo da garantire continuità su reti con throughput limitato. Successivamente si
codificano altri flussi detti enhancement che aggiungono progressivamente sempre più informazioni di perfezionamento del
flusso base per aumentare frame rate, risoluzione e qualità video.

Il server trasmette contemporaneamente sia il flusso base che i flussi enhancement. Se si dovessero verificare fenomeni
di congestione della rete, il server o i nodi intermedi possono decidere di interrompere la trasmissione di uno o più flussi
enhancement. In questo modo si abbassa il bitrate e si mantiene la continuità di trasmissione del flusso base per garantire
fluidità del video.

\subsubsection*{Adaptive bitrate - ABR}
La soluzione ad oggi più adottata è quella del client-driven adaptive bitrate (ABR). Consiste in un protocollo di
``collaborazione'' a livello applicativo tra client e server senza richiedere requisiti all'infrastruttura di rete
e ai livelli inferiori.

In pratica il server suddivide il video in segmenti di durata fissa, tipicamente 2-10 secondi, e li codifica in più varianti
con bitrate e qualità differenti e li offre al client. L'applicativo dal lato client si occupa di decidere quale versione
scaricare in base alle prestazioni della rete e alla qualità video desiderata.

Per permettere al client di effettuare la scelta del segmento migliore da scaricare, il server mette a disposizione
un file detto Media Description Presentation (MDP) con le informazioni necessarie, come:
\begin{itemize}
	\item numero di segmenti disponibili e durata di ciascun segmento
	\item numero di livelli di qualità disponibili per ciascun segmento
	\item bitrate, risoluzione e frame rate di ciascun livello
	\item l'url di ciascun segmento per ciascun livello
\end{itemize}

\subsection{Dynamic adaptive streaming over HTTP - DASH}
\subsubsection*{Descrizione del protocollo}
L'implementazione più diffusa dell'approccio ABR è il protocollo DASH, definito nello standard MPEG-DASH. Il protocollo
descrive un modello pull, dove il client inizialmente richiede al server il file MDP e progressivamente scarica dal server
i segmenti video che ritiene migliori uno alla volta. La scelta del segmento viene fatta lato client. I dati viaggiano
su protocolli di trasporto HTTP/3 e QUIC, secondo il modello pull (il client richiede i segmenti al server).

\subsubsection*{Vantaggi architetturali di HTTP/3 e QUIC rispetto a RTSP/RTMP e TCP}
L'utilizzo di HTTP/3 e QUIC (basati su UDP) permette di superare le limitazioni dello streaming su protocolli legacy RTSP/RTMP
e del trasporto basato su TCP, introducendo i seguenti miglioramenti:
\begin{itemize}
	\item \textbf{approccio web-centrico}: si abbandona il modello push e stateful su porte non standard web dei formati
	RTSP/RTMP, in favore di un'architettura pull e stateless in cui la gestione del traffico è affidata al client; l'uso
	di porte standard (80/443) garantisce l'attraversamento trasparente di NAT e firewall, e permette la simbiosi con le
	CDN per il caching
	\item \textbf{assenza di HOL blocking}: il multiplexing nativo di UDP permette di avere più flussi indipendenti e la
	perdita di un pacchetto non blocca l'intera coda di trasmissione come nel TCP
	\item \textbf{assenza di latenza di Setup}: vengono eliminati i multipli handshake sequenziali imposti da TCP,
	consentendo l'avvio quasi istantaneo della trasmissione dati
	\item \textbf{connection migration}: la connessione non è vincolata alla coppia IP/Porta, garantendo la continuità
	dello streaming anche durante i cambi di rete (es. transizione da Wi-Fi a rete mobile).
	\item \textbf{evoluzione rapida:} lo spostamento della logica di controllo della congestione dal kernel allo strato
	applicativo facilita i futuri aggiornamenti del protocollo
\end{itemize}

\subsubsection*{Stato di salute del buffer e gestione del rebuffering}
Si definisce lo stato di salute del buffer del client \(B(t)\) come la quantità di secondi di video presenti nel buffer
all'istante \(t\). Quando \(B(t) = 0\) il buffer è vuoto ed è necessario effettuare un rebuffering, ovvero interrompere
la riproduzione per riempire il buffer con un certo numero (\(L\) o \(M\)) di nuovi segmenti.
\begin{itemize}
	\item \(L\): \# segmenti da scaricare durante il primo rebuffering, prima di iniziare la riproduzione
	\item \(M\): \# segmenti da scaricare durante un rebuffering intermedio, prima di riprendere la riproduzione
\end{itemize}
I parametri \(M\) e \(L\) vengono scelti dal servizio di streaming. Più grandi sono, maggiore è la durata del rebuffering,
ma minore sarà la probabilità di un nuovo rebuffering successivo.

In condizione di playout, ovvero durante la riproduzione del video, lo stato del buffer \(B(t)\) decresce al tasso di
un secondo al secondo a cui viene sommata la quantità di video scaricato. In condizioni di rebuffering, invece, lo stato
del buffer \(B(t)\) aumenta al tasso di video scaricato.
\[\frac{d \, B_\text{playout}(t)}{d \, t} = \frac{S}{R_C} - 1 \quad \text{con} \; \begin{array}{c}
	S < R_C \rightarrow \text{\small buffer in svuotamento} \\
	S > R_C \rightarrow \text{\small buffer in riempimento}
\end{array} \quad\;\;\Big|\;\;\quad \frac{d \, B_\text{rebuffering}(t)}{d \, t} = \frac{S}{R_C}\]

\subsubsection*{Tempo di download di un segmento}
Si definisce il tempo di download \(T_D(n,k)\) del segmento \(n\) al livello di qualità \(k\) come rapporto tra il tempo
di download (prodotto tra la durata del segmento in secondi \(T_S\) e il tasso di codifica \(R_C(k)\)) e il throughput della
rete \(S_n\) durante il download del segmento \(n\):
\[T_D(n,k) = \frac{T_S \cdot R_C(k)}{S_n} \;\; \implies \;\; \frac{T_D(n,k)}{T_S} = \frac{R_C(k)}{S_n}\]

\subsubsection*{Quality of Experience - QoE}
La Quality of Experience (QoE) è una metrica che misura la qualità percepita dall'utente finale durante la fruizione
di un contenuto multimediale ed è data dalla seguente formula:
\[J = \underbrace{\; \sum_{n=1}^{N} q(n) \;}_\text{qualità del video} - \underbrace{\; \lambda \cdot \sum_{n=1}^{N} \left| q(n) - q(n-1) \right| \;}_\text{penalizzazione per cambio di qualità} - \underbrace{\; \mu \cdot \sum_{n=1}^{N} T_\text{stall}(n) \;}_\text{penalizzazione per gli stalli} \]
\begin{itemize}
	\item \(q(n)\): qualità del segmento \(n\) in base al livello di codifica scelto
	\item \(\lambda\): peso della penalizzazione per il cambio di qualità tra due segmenti consecutivi
	\item \(\mu\): peso della penalizzazione per gli stalli dovuti al rebuffering (solitamente molto alto)	
\end{itemize}
